\section{Retargeting and Character Diversity}

For the first demo, I used a single walking motion capture clip and transferred it onto three different MetaHumans with noticeably different proportions. The three characters were my own avatar, built as a roughly six-foot, slightly chubby male and presented using Epic Games' MetaHuman Techwear Outfit \cite{fab_avatar_outfit}, together with two alternative MetaHumans: a short female character using the Casual Hoodie \& Jeans outfit by polycornStudio \cite{fab_female_outfit}, and an older male character using the Poor Man Outfit by NOETA Marketstore \cite{fab_oldman_outfit}. Since the source clip came from my own mocap recording, it was naturally best suited to my own avatar, which made it a useful baseline for seeing what changed when the same motion was applied to characters with very different body types. That also pushed this away from being just a transfer test and toward what had to happen afterwards, because what looks acceptable on one character started to feel too light, too stretched, or badly grounded on another.

\subsection{Retargeting Workflow}

The initial workflow followed Unreal Engine's standard retargeting pipeline \cite{epic_ik_retargeter}. The source mocap skeleton was imported, IK rigs were generated for both source and target characters, and an IK Retargeter was created to transfer the motion. At a basic level this worked, but the first results exposed several problems:

\begin{itemize}[leftmargin=*]
    \item root misalignment that caused the character to disappear or collapse visually,
    \item pelvis height errors that pushed the body too low,
    \item foot penetration into the ground,
    \item arm offsets caused by pose mismatch between source and target skeletons,
    \item and instability introduced by IK goals and solver settings that did not suit the source rig.
\end{itemize}

These issues made it clear that the default setup was not robust enough for this dataset, so a debugging pass was needed before the retargeting results became usable. IK goals that referenced incompatible bones were removed, the Full Body IK solver was disabled in the retargeting configuration, and the retarget root was changed to the pelvis or hips. I also found that Unreal's post-process animation logic could interfere with testing, so that system had to be disabled temporarily while verifying the raw retarget result.

\begin{figure}[H]
    \centering
    \includegraphics[width=0.82\linewidth]{Demo1_Retarget.png}
    \caption{IK Retargeter pose transfer from mocap to MetaHumans.}
\end{figure}

\subsection{Per-Character Adaptation}

My initial assumption had been that I could retarget the animation once and then reuse the result across all characters. In practice, this produced uneven quality because the same source walk was interpreted very differently depending on the target character's build. My own avatar needed the least correction overall, since the mocap came from my own performance, while the shorter female character and the older male character each revealed much clearer exaggerations in posture, stride length, and grounding.

Because of this, I changed the workflow and created a separate retargeting pass for each MetaHuman rather than relying on a single shared export. This allowed me to make small but important adjustments per character, including pelvis height correction, posture changes, root alignment tweaks, and later playback adjustments. In practice, my own avatar stayed closest to the original clip and therefore kept its original speed, the female character was made to feel lighter and more energetic in the stride but ended up only very slightly slower because of the shorter step length, and the older male character was clearly the slowest, with stiffer steps and a more obviously adjusted posture. Although this increased workload, it gave far more control and produced noticeably cleaner results.

\subsection{Sequencer and Control Rig Editing}

The retargeted clips were then taken into Sequencer and baked to Control Rig for cleanup using the MetaHuman Control Rig workflow \cite{epic_metahuman_control_rig}. The main corrections during this stage were foot placement adjustments to reduce ground penetration and sliding, fixes to the hips, spine, and shoulder rotations, visible jitter cleanup, posture refinement for each character, and smoothing of awkward moments in the stride cycle. These edits improved grounding and made the differences between the characters more apparent instead of all three looking like simple copies of the same clip.

\begin{figure}[H]
    \centering
    \begin{subfigure}[t]{0.48\linewidth}
        \centering
        \includegraphics[width=\linewidth]{Demo1_Posture.png}
        \caption{Posture cleanup.}
    \end{subfigure}
    \hfill
    \begin{subfigure}[t]{0.48\linewidth}
        \centering
        \includegraphics[width=\linewidth]{Demo1_Stride.png}
        \caption{Stride distance and grounding.}
    \end{subfigure}
    \caption{Sequencer keyframe adjustments for posture, stride distance, and speed.}
\end{figure}

To strengthen the distinction between characters, I also adjusted play rate and stance slightly per character. My own avatar stayed closest to the original clip, the shorter female character was shaped to feel lighter and more energetic even though the shorter step length made her read marginally slower overall, and the older male character was pushed toward a slower and stiffer result. 

\subsection{Critical Evaluation}

The most useful thing this demo showed is that retargeting is not always simple. The transfer itself worked, but the raw result still was not strong enough to present as finished motion, and the biggest lesson here was how strongly body proportions affect perception even when the underlying source clip stays exactly the same. It also shows that once proportions start changing, the motion also needs to be adapted, and the combination of per-character retargeting and Control Rig cleanup made a much bigger difference than the transfer alone.

The weakness is that the workflow is still very manual, because different characters produce different grounding and pose issues and the amount of cleanup scales up quickly once more than one body type is involved. Even after cleanup, there can still be small problems like slight sliding or imperfect foot contact.

Overall, I think the demo works because it shows both the usefulness and the limits of retargeting honestly. It makes the case that believable reused motion usually comes from retargeting plus editing, not retargeting on its own.
